\documentclass[11pt,a4paper]{article}
\usepackage[T1]{fontenc}
\usepackage[utf8]{inputenc}
\usepackage[french]{babel}
\usepackage{lmodern}
\usepackage{geometry}
\usepackage{longtable}
\usepackage{array}
\usepackage{hyperref}
\usepackage{amsmath}
\geometry{margin=2.5cm}

\title{Inventaires fongiques -- Sp\'ecification fonctionnelle (ICR)}
\author{Eddy Boite}
\date{9 juillet 2026 -- v1.4}

\begin{document}
\maketitle

\section*{Objectif}
Ce document formalise les exigences fonctionnelles du script \texttt{scripts/Inventaires\_completude\_representativite.R} (v1.4) pour l'analyse de compl\'etude et de repr\'esentativit\'e d'inventaires fongiques.

Fonctions livr\'ees :
\begin{enumerate}
  \item courbe temps-esp\`eces,
  \item ajustements lin\'eaire/hyperbolique (si conditions remplies),
  \item estimation d'asymptote,
  \item calcul de \texttt{TEE} et de $Ir = 1 - TEE$,
  \item fr\'equences d'occurrence par esp\`ece,
  \item occupation spatiale (si \texttt{placette}),
  \item CA/AFC (si \texttt{vegan}),
  \item m\'etriques de pertinence scientifique,
  \item manifeste final de production.
\end{enumerate}

\section*{Entr\'ees}
\subsection*{Fichier principal}
\texttt{data/observations.csv} (ou \texttt{.txt}/\texttt{.tsv})

\subsection*{Colonnes}
\begin{itemize}
  \item Obligatoires : \texttt{site}, \texttt{date}, \texttt{visite\_id}, \texttt{espece}
  \item Optionnelle : \texttt{placette}
\end{itemize}

\subsection*{R\`egles fonctionnelles}
\begin{itemize}
  \item Parsing date via \texttt{INVENTAIRES\_DATE\_FORMAT} (d\'efaut : \texttt{\%Y-\%m-\%d}).
  \item D\'eduplication sur \texttt{(site, placette, date, visite\_id, espece)}.
  \item Visite distincte : \texttt{date + visite\_id} (et \texttt{site} en multi-sites).
\end{itemize}

\section*{Configuration centralis\'ee}
Tous les param\`etres du pipeline sont d\'efinis dans la fonction \texttt{get\_embedded\_config()} et affich\'es dans l'en-t\^ete du script au d\'emarrage. La configuration supporte les surcharges via variables d'environnement \texttt{INVENTAIRES\_*} :

\begin{itemize}
  \item \texttt{INVENTAIRES\_INPUT\_FILE} : chemin du fichier source
  \item \texttt{INVENTAIRES\_OUTPUT\_DIR} : r\'epertoire de sortie
  \item \texttt{INVENTAIRES\_DATE\_FORMAT} : format des dates
  \item \texttt{INVENTAIRES\_CSV\_STRICT} : mode strict (booléen)
  \item \texttt{INVENTAIRES\_CSV\_ALLOW\_EXTRA\_COLS} : tolérance colonnes suppl\'ementaires
  \item \texttt{INVENTAIRES\_FREQ\_BREAKS} : s\'equence des intervalles de fr\'equence
  \item \texttt{INVENTAIRES\_FREQ\_LABELS} : \'etiquettes des intervalles
  \item \texttt{INVENTAIRES\_MAKE\_CA} : activer CA/AFC (booléen)
  \item \texttt{INVENTAIRES\_MIN\_VISITS\_FOR\_MODEL} : seuil de mod\'elisation
  \item \texttt{INVENTAIRES\_WIDTH}, \texttt{INVENTAIRES\_HEIGHT}, \texttt{INVENTAIRES\_DPI} : param\`etres graphiques
\end{itemize}

Aucune modification du code source n'est n\'ecessaire : les param\`etres peuvent \^etre adjust\'es par variable d'environnement.

\section*{Validation CSV}
Le contr\^ole de conformit\'e est ex\'ecut\'e avant traitement.
\begin{itemize}
  \item \texttt{results/ICR/ICR\_00\_csv\_conformite\_report.csv} (toujours)
  \item \texttt{results/ICR/ICR\_00\_csv\_conformite\_problems.csv} (si anomalies)
\end{itemize}
En mode strict (\texttt{INVENTAIRES\_CSV\_STRICT=TRUE}), la non-conformit\'e est bloquante.

\section*{Sorties attendues}
\subsection*{Organisation thématique}
Les résultats sont organisés selon une hiérarchie thématique sous \texttt{results/ICR/} :

\begin{itemize}
  \item \texttt{00\_data\_quality/} : rapports d'audit CSV et conformité
  \item \texttt{01\_richesse/}, \texttt{02\_completude\_repr/}, \texttt{03\_frequences/} : analyses par site (réservées)
  \item \texttt{04\_spatial/}, \texttt{05\_temporal/} : analyses spéciales (réservées)
  \item \texttt{06\_metrics/} : métriques scientifiques globales
  \item \texttt{comparisons/} : visualisations comparatives multi-sites
  \item \texttt{sites/<site\_name>/} : **TOUS les fichiers par site** (organisés en un seul répertoire)
\end{itemize}

\subsection*{Globales (\texttt{results/ICR/})}
\begin{itemize}
  \item \texttt{ICR\_donnees\_preparees.csv}
  \item \texttt{ICR\_resume\_tous\_sites.csv}
  \item \texttt{ICR\_metrics\_pertinence\_tous\_sites.csv}
  \item \texttt{ICR\_comparaison\_completude\_sites.png} (dans \texttt{comparisons/})
  \item \texttt{ICR\_comparaison\_ir\_sites.png} (dans \texttt{comparisons/})
  \item \texttt{ICR\_metrics\_pertinence\_heatmap\_sites.png} (dans \texttt{comparisons/})
  \item \texttt{ICR\_metrics\_pertinence\_score\_sites.png} (dans \texttt{comparisons/})
\end{itemize}

\subsection*{Par site (\texttt{results/ICR/sites/<site\_name>/})}
Tous les fichiers CSV et PNG d'un site sont consolidés dans un seul répertoire :
\begin{itemize}
  \item \texttt{ICR\_01\_courbe\_temps\_especes.csv} et graphique PNG
  \item \texttt{ICR\_02\_tee\_ir.csv} et graphique PNG
  \item \texttt{ICR\_03\_frequence\_especes.csv} et graphique PNG
  \item \texttt{ICR\_04\_resume\_site.csv}
  \item \texttt{ICR\_07\_metrics\_pertinence.csv} et graphique PNG
  \item \texttt{ICR\_05\_modeles.csv} (si conditions remplies)
  \item \texttt{ICR\_06\_occupation\_spatiale.csv} (si placette)
  \item \texttt{ICR\_05\_temporel\_vs\_spatial.png} (conditionnelle)
  \item \texttt{ICR\_06\_CA\_placettes\_especes.png} (si vegan et placette)
\end{itemize}

\section*{Indicateurs}
\subsection*{TEE}
Proportion d'esp\`eces observ\'ees une seule fois.

\subsection*{Indice de repr\'esentativit\'e}
\[
Ir = 1 - TEE
\]
Rep\`eres : \texttt{Ir < 0.60} (faible), \texttt{0.60 <= Ir < 0.80} (moyenne), \texttt{Ir >= 0.80} (bonne).

\subsection*{Compl\'etude}
\[
\text{Compl\'etude} = \frac{S_{obs}}{S_{asymptote}}
\]
Rep\`eres : \texttt{< 0.70} (insuffisante), \texttt{0.70--0.90} (interm\'ediaire), \texttt{>= 0.90} (avanc\'ee).

\section*{Logging syst\'ematis\'e}
Chaque ex\'ecution du script g\'en\'ere un fichier log horodat\'e dans le r\'epertoire \texttt{logs/} du projet :

\begin{itemize}
  \item \textbf{Format du nom} : \texttt{ICR\_YYYYMMDD\_HHMM.log}
  \item \textbf{Format des ent\'ees} : \texttt{[YYYY-MM-DD HH:MM:SS] [LEVEL] message}
  \item \textbf{Niveaux} : \texttt{INFO}, \texttt{DEBUG}, \texttt{WARN}
\end{itemize}

Les logs enregistrent :
\begin{enumerate}
  \item En-t\^ete du script (nom, version, timestamp, répertoire de travail, version R)
  \item Chaque étape du pipeline (résolution fichier, audit CSV, analyses par site)
  \item Les parameters de configuration actifs au démarrage
  \item Le manifeste final des livrables avec statuts (\texttt{OK}, \texttt{MANQUANT}, \texttt{OPTIONNEL\_NON\_GENERE})
\end{enumerate}

Les avertissements sont également affich\'es en console et dans les logs en niveau \texttt{WARN}.

\section*{Variables d'environnement support\'ees}
\begin{itemize}
  \item \texttt{INVENTAIRES\_INPUT\_FILE}
  \item \texttt{INVENTAIRES\_OUTPUT\_DIR}
  \item \texttt{INVENTAIRES\_DATE\_FORMAT}
  \item \texttt{INVENTAIRES\_CSV\_STRICT}
  \item \texttt{INVENTAIRES\_CSV\_ALLOW\_EXTRA\_COLS}
  \item \texttt{INVENTAIRES\_AUTO\_RUN}
\end{itemize}

\section*{Modes d'ex\'ecution}
Le script peut \^etre ex\'ecut\'e en mode standard (recommand\'e) ou avec des flags internes :
\begin{itemize}
  \item \texttt{DEBUG\_MODE = TRUE} : verbosit\'e accrue et affichage d\'etaill\'e des avertissements.
  \item \texttt{BENCHMARK\_MODE = TRUE} : mesure des temps sur les calculs principaux.
  \item \texttt{CLEAN\_ENVIRONMENT = TRUE} : nettoyage de l'environnement avant ex\'ecution (usage expert).
\end{itemize}

Le comportement d'auto-lancement est pilot\'e par \texttt{INVENTAIRES\_AUTO\_RUN} (d\'efaut : \texttt{TRUE}).

\section*{Am\'eliorations du script (v1.4)}
\subsection*{Phase 1 : Fiabilisation}
Gestion robuste des formats de fichiers hétérogènes, déduplication intelligente des observations, traitement des dates flexibles, gestion des erreurs CSV granulaire.

\subsection*{Phase 2 : Performance}
Optimisation du modèle hyperbolique via \texttt{minpack.lm::nlsLM}, calculs vectorisés par site, graphiques pré-calculés en mémoire avant export PNG.

\subsection*{Phase 3 : Graphiques et visualisations}
Cartes de chaleur multi-sites, tableaux de bord par site (7 graphiques), comparaisons visuelles de complétude, analyses de CA/AFC si \texttt{vegan} disponible.

\subsection*{Phase 4 : Industrialisation}
\begin{enumerate}
  \item \textbf{Configuration centralisée} : fonction \texttt{get\_embedded\_config()} avec 14+ paramètres, surchargeable via variables d'environnement
  \item \textbf{Logging systématisé} : fichiers logs horodatés avec métadonnées d'exécution, manifeste final des livrables
  \item \textbf{Organisation thématique} : hiérarchie logique /sites/, /comparisons/, répertoires réservés pour extensions futures
  \item \textbf{Métriques scientifiques} : pertinence des espèces, scores de complétude détaillés, dashboards analytiques
\end{enumerate}

\section*{Journalisation fonctionnelle}
La sortie est journalis\'ee en console (niveaux info/debug/warning), puis un manifeste final indique l'\'etat des livrables :
\texttt{OK}, \texttt{MANQUANT}, \texttt{OPTIONNEL\_NON\_GENERE}.

Le script applique enfin une politique stricte de nettoyage du fichier parasite \texttt{Rplots.pdf} lorsqu'il est g\'en\'er\'e pendant l'ex\'ecution.

\section*{D\'ependances}
Requises : \texttt{dplyr}, \texttt{tidyr}, \texttt{ggplot2}, \texttt{purrr}, \texttt{readr}, \texttt{stringr}, \texttt{forcats}, \texttt{tibble}, \texttt{scales}, \texttt{gridExtra}.\\
Optionnelles : \texttt{minpack.lm} (mod\`ele hyperbolique), \texttt{vegan} (CA/AFC).

\end{document}
